\documentclass{beamer}
\usepackage[utf8]{inputenc}
\usetheme{AMU}
\usepackage{amsmath,amssymb}
\usepackage{array}
\usepackage{booktabs}
\usepackage{tikz}
\usetikzlibrary{calc,positioning}

\setbeamertemplate{footline}{%
  \leavevmode%
  \hbox to .88\paperwidth{%
    \hspace*{1em}{\color{AMU_dk2}\footnotesize\insertshorttitle}%
    \hfil%
    {\color{AMU_dk2}\footnotesize\insertframenumber{} / \inserttotalframenumber}%
  }%
  \vskip1ex%
}

\title[Stochastic Models and Simulation]{Stochastic Models and Simulation}
\subtitle{Week 4: how long does an agent stay busy?}
\author{Sebastian M\"uller}
\date{2026/27}

\begin{document}

\begin{frame}
  \renewcommand{\AMUtitlebar}{%
    \noindent\colorbox{AMU_dk2}{\rule{0pt}{1.2ex}\hspace{\dimexpr0.36\linewidth-2\fboxsep\relax}}%
    \colorbox{AMU_lt2}{\rule{0pt}{1.2ex}\hspace{\dimexpr0.28\linewidth-2\fboxsep\relax}}%
    \colorbox{AMU_acc6}{\rule{0pt}{1.2ex}\hspace{\dimexpr0.36\linewidth-2\fboxsep\relax}}%
  }
  \titlepage
\end{frame}

\section{Observe service}

\begin{frame}{Does the six-minute assumption hold up?}
  In Week 3, every simulated service duration came from an exponential
  distribution with a mean of six minutes.

  \vspace{0.7em}
  \begin{block}{Why look at service times?}
    While an agent handles a long call, the next callers have one fewer
    agent available. Both the typical duration and the frequency of long
    calls can change the staffing comparison.
  \end{block}

  \vspace{0.4em}
  Today we measure service, fit two distributions, and check which calls
  account for the most work. Then we revisit the staffing experiment.
\end{frame}

\begin{frame}{Three inputs to the bank queue}
  \centering
  \begin{tikzpicture}[
    inputbox/.style={draw=AMU_dk2,rounded corners,align=center,
      text width=2.4cm,minimum height=1.25cm,fill=AMU_dk2!7}]
    \node[inputbox,minimum height=1.7cm] (arrival) at (-3.2,1.7) {Agent requests\\When do they arrive?};
    \node[inputbox,minimum height=1.7cm,fill=AMU_lt2,draw=AMU_acc6,very thick] (service) at (0,1.7)
      {Service duration\\\textbf{Week 4}};
    \node[inputbox,minimum height=1.7cm] (patience) at (3.2,1.7) {Caller patience\\When do callers leave?};
    \node[inputbox,text width=6.3cm] (queue) at (0,-0.55)
      {Queue and staffing\\Predict waiting and abandonment};
    \draw[->,thick] (arrival.south) -- ([xshift=-2.45cm]queue.north);
    \draw[->,thick] (service.south) -- (queue.north);
    \draw[->,thick] (patience.south) -- ([xshift=2.45cm]queue.north);
  \end{tikzpicture}

  \vspace{0.8em}
  \raggedright
  Next week we count actual agent requests and study how long callers are
  willing to wait before hanging up.
\end{frame}

\begin{frame}{Which part of the call are we measuring?}
  \begin{center}
    \begin{tabular}{@{}lll@{}}
      \toprule
      Call entered queue & Agent started & Agent finished\\
      \midrule
      09:14:12 & 09:14:54 & 09:18:24\\
      \bottomrule
    \end{tabular}
  \end{center}
  \small Example with invented times. Outcome \texttt{AGENT} means a human
  agent handled the call.\normalsize

  \vspace{0.65em}
  \begin{itemize}
    \item The caller waited \textbf{42 seconds}.
    \item The agent was busy for \textbf{210 seconds}. This is the recorded service duration,
      \texttt{ser\_time}.
    \item If the caller leaves before service, there is no agent service duration.
  \end{itemize}
\end{frame}

\begin{frame}{Same mean service time, different waits}
  One agent. Three calls arrive at minutes $0$, $4$, and $8$.
  Both schedules contain 12 minutes of service.

  \vspace{0.7em}
  \centering
  \begin{tikzpicture}[x=0.56cm,y=0.88cm,font=\small]
    \node[below,align=center] at (6,3.3) {Call arrivals};
    \foreach \x in {0,4,8} {
      \draw[->,thick] (\x,2.8) -- (\x,2.2);
      \draw[densely dotted,black!45] (\x,2.1) -- (\x,0.12);
    }
    \node[left] at (-0.3,1.68) {4, 4, 4 min};
    \node[left] at (-0.3,0.70) {8, 2, 2 min};
    \foreach \x/\finish/\caller in {0/4/1,4/8/2,8/12/3} {
      \draw[fill=AMU_dk2!16,draw=AMU_dk2] (\x,1.4) rectangle
        node {\caller} (\finish,1.95);
    }
    \foreach \x/\finish/\caller in {0/8/1,8/10/2,10/12/3} {
      \draw[fill=AMU_lt2,draw=AMU_acc6] (\x,0.42) rectangle
        node {\caller} (\finish,0.97);
    }
    \draw[->] (0,0) -- (12.55,0) node[right] {min};
    \foreach \x in {0,4,8,10,12} {
      \draw (\x,-0.06) -- (\x,0.06);
      \node[below] at (\x,-0.06) {\x};
    }
  \end{tikzpicture}

  \vspace{0.5em}
  \raggedright
  Each block shows when one caller is in service. How long does each
  caller wait in the second schedule? What changed for callers 2 and 3?
\end{frame}

\begin{frame}{Check the very short services before fitting}
  \small
  Compare \texttt{AGENT} calls on Sunday--Thursday whose service started
  from 09:00 to before 11:00. Keep the full service duration, even if
  service ends after 11:00.

  \vspace{0.55em}
  \begin{center}
    \begin{tabular}{@{}lrrr@{}}
      \toprule
      Month (1999) & Services & Under 10 s & Share\\
      \midrule
      January  & 3,631 & 163 & 4.49\%\\
      November & 5,072 &  28 & 0.55\%\\
      \bottomrule
    \end{tabular}
  \end{center}

  Brown et al.\ (2005) analysed this same bank. They report that its
  manager traced many short services to agents hanging up.

  \vspace{0.45em}
  For our fit, exclude services under 10 seconds. This can also remove
  genuine brief calls. In November, we retain 2,530 services on 1--14
  (fit) and 2,514 on 15--30 (check).

  \vspace{0.5em}
  {\scriptsize
    \href{https://seelab.net.technion.ac.il/data/}{Technion SEE Lab, AnonymousBank}.
    \href{https://faculty.wharton.upenn.edu/wp-content/uploads/2012/04/Brown-etal-CCdata.pdf}{Brown et al.\ (2005), \emph{Statistical Analysis of a Telephone Call Center}, \S4.1}.}
\end{frame}

\begin{frame}{Most services are short. Some run much longer.}
  \centering
  \begin{tikzpicture}[x=1.45cm,y=0.0040cm]
    \node[anchor=west,font=\scriptsize] at (-0.6,870) {Service records};
    \foreach \y in {0,200,400,600,800} {
      \draw[black!18] (-0.62,\y) -- (5.6,\y);
      \node[anchor=east,font=\scriptsize] at (-0.62,\y) {\y};
    }
    \foreach \x/\n/\lab in {0/423/10 s--1 min,1/764/1--2 min,2/745/2--4 min,
                             3/254/4--6 min,4/203/6--10 min,5/141/10+ min} {
      \fill[AMU_dk2!78] (\x-0.29,0) rectangle (\x+0.29,\n);
      \node[above,font=\scriptsize] at (\x,\n) {\n};
      \node[below,font=\scriptsize] at (\x,0) {\lab};
    }
  \end{tikzpicture}

  \vspace{0.55em}
  \raggedright\small
  Counts of the 2,530 retained services on training dates. The bars show
  counts in ranges of different widths. Each range includes its lower endpoint.
\end{frame}

\section{Fit and check}

\begin{frame}{Fit an exponential baseline to the observed mean}
  Let $S$ be the service duration in seconds. Fit to the retained
  records from 1--14 November.

  \vspace{0.65em}
  \begin{itemize}
    \item Fit an exponential distribution with the observed mean,
      $205.7$ seconds, or about $3.43$ minutes.
    \item An exponential has equal mean and standard deviation. The
      observed standard deviation is larger, at $260.2$ seconds.
  \end{itemize}

\end{frame}

\begin{frame}{Give the model a separate control for spread}
  \small
  Draw $Z$ from a normal distribution with mean $\mu$ and standard deviation
  $\sigma$. Set $S=e^Z$ seconds. Then the service duration $S$ is lognormal.

  \vspace{0.5em}
  Let $m$ and $v$ be the empirical mean and variance of $S$ on the
  training dates. Pick the two parameters to match those two quantities:
  \[
    \sigma^2=\log(1+v/m^2), \qquad \mu=\log m-\sigma^2/2.
  \]

  The exponential and lognormal candidates now have the same average
  duration. They can predict different shares of long calls.

  \vspace{0.45em}
  {\scriptsize
  Brown et al.\ found a close lognormal approximation for November--December.
  Here we fit our own parameters for the selected morning services and test them
  on later November dates.}
\end{frame}

\begin{frame}{How often does a service last at least ten minutes?}
  \small
  We fit both models on 1--14 November. We set aside the services from
  15--30 November to check the predictions.

  \vspace{0.65em}
  \centering
  \begin{tabular}{@{}rrrrr@{}}
    \toprule
    & \multicolumn{2}{c}{Observed in November 1999}
    & \multicolumn{2}{c}{Model predictions}\\
    \cmidrule(lr){2-3}\cmidrule(l){4-5}
    At least & 1--14 (fit) & 15--30 (check) & Exponential & Lognormal\\
    \midrule
    2 min  & 53.1\% & 59.9\% & 55.8\% & 52.5\%\\
    6 min  & 13.6\% & 19.1\% & 17.4\% & 14.4\%\\
    \textbf{10 min} & \textbf{5.6\%} & \textbf{8.3\%} & \textbf{5.4\%} & \textbf{5.7\%}\\
    15 min &  2.1\% &  3.4\% &  1.3\% &  2.3\%\\
    \bottomrule
  \end{tabular}

  \vspace{0.65em}
  \raggedright
  On 15--30 November, 208 of 2,514 retained services lasted at least ten
  minutes (8.3\%). Both models predict a lower share.
\end{frame}

\begin{frame}{8.3\% of calls account for 35.8\% of service time}
  Here we call a service \textbf{long} if it lasts \textbf{10 minutes or
  more}. Both bars refer to the retained calls from 15--30 November.

  \vspace{0.6em}
  \centering
  \begin{tikzpicture}[x=0.064cm,y=0.95cm,font=\small]
    \node[anchor=east] at (-3,1.3) {By call count};
    \node[anchor=east] at (-3,0.4) {By service hours};
    \draw[fill=AMU_dk2!7,draw=black!24] (0,1.07) rectangle (100,1.53);
    \draw[fill=AMU_dk2!7,draw=black!24] (0,0.17) rectangle (100,0.63);
    \fill[AMU_dk2!82] (0,1.07) rectangle (8.27,1.53);
    \fill[AMU_dk2!82] (0,0.17) rectangle (35.76,0.63);
    \node[right] at (8.27,1.3) {\textbf{8.3\%}};
    \node[right] at (35.76,0.4) {\textbf{35.8\%}};
    \node[below] at (0,0.17) {0\%};
    \node[below] at (100,0.17) {100\%};
  \end{tikzpicture}

  \vspace{0.6em}
  \raggedright
  There were 208 long calls out of 2,514. Those same 208 calls took
  about 63 of the 176 total service hours recorded for these calls.
\end{frame}

\begin{frame}{Did the mix of services change?}
  \small
  The bank records six service types. We combine them into two groups
  for a first comparison:

  \begin{itemize}
    \item \textbf{Regular banking:} standard banking (\texttt{PS}) and banking
      in English (\texttt{PE}).
    \item \textbf{Other services:} internet help (\texttt{IN}), stock exchange
      (\texttt{NE}), prospective customers (\texttt{NW}), and callbacks put
      on hold (\texttt{TT}).
  \end{itemize}

  \vspace{0.25em}
  \centering
  \begin{tabular}{@{}llrr@{}}
    \toprule
    Dates (Nov.) & Service group & Share of calls & Long within group\\
    \midrule
    1--14 & Regular banking & 70.4\% & 4.2\%\\
          & Other services  & 29.6\% & 8.9\%\\[0.2em]
    15--30 & Regular banking & 68.3\% & 6.7\%\\
           & Other services  & 31.7\% & 11.7\%\\
    \bottomrule
  \end{tabular}

  \vspace{0.55em}
  \raggedright
  Could the change in the mix alone explain the rise in long services?
\end{frame}

\begin{frame}{Can the service mix explain the rise in long calls?}
  \small
  Use the mix from \textbf{15--30 November}: 68.3\% regular banking
  and 31.7\% other services. Keep the share of long calls \textbf{within
  each group} at its 1--14 November value.

  \[
    \underbrace{0.683 \times 4.2\%}_{\text{Regular banking}}
    + \underbrace{0.317 \times 8.9\%}_{\text{Other services}}
    \approx 5.7\%.
  \]

  \vspace{0.15em}
  \centering
  \begin{tabular}{@{}lr@{}}
    \toprule
    & Share lasting at least 10 min\\
    \midrule
    Observed on 1--14 November & 5.6\%\\
    If only the mix changed    & 5.7\%\\
    Observed on 15--30 November & 8.3\%\\
    \bottomrule
  \end{tabular}

  \vspace{0.65em}
  \raggedright
  Changing the mix alone adds about 0.1 percentage points. The observed
  rise was 2.7 points, and long calls became more common within both groups.
\end{frame}

\section{Simulate and decide}

\begin{frame}{Pair the simulations, one draw per caller}
  \small
  Generate one morning of arrivals as in Week 3. Draw a fresh uniform
  number between 0 and 1 for \textbf{each caller}. Each model converts
  that draw to a duration.

  \vspace{0.5em}
  Here are three example callers. A draw of 0.95 picks the duration
  by which 95\% of services finish under each model.

  \vspace{0.5em}
  \centering
  \begin{tabular}{@{}lrrrr@{}}
    \toprule
    & & \multicolumn{3}{c}{Service duration (minutes)}\\
    \cmidrule(l){3-5}
    Caller & Draw & Week 3 exp. & Fitted exp. & Fitted lognormal\\
    \midrule
    1 & 0.25 & 1.73 & 0.99 & 1.10\\
    2 & 0.60 & 5.50 & 3.14 & 2.72\\
    3 & 0.95 & 17.97 & 10.27 & 10.61\\
    \bottomrule
  \end{tabular}

  \vspace{0.6em}
  \raggedright
  \textbf{Recall Week 3:} pairing helps reduce random variation in the
  comparison. Across models, share the arrivals and each caller's draw.
  Within a model, share the arrivals and service durations at 16 and
  17 agents. Repeat for 200 mornings.
\end{frame}

\begin{frame}{How does the staffing comparison change?}
  \small
  Compute one mean wait per morning for each staffing level. The table
  averages those 200 morning means and their paired differences, in minutes.

  \vspace{0.7em}
  \centering
  \begin{tabular}{@{}lrrrl@{}}
    \toprule
    Service input & Wait 16 & Wait 17 & Gain & 95\% MC interval\\
    \midrule
    Week 3 exponential & 4.354 & 2.749 & 1.604 & [1.413, 1.796]\\
    Fitted exponential & 0.023 & 0.010 & 0.013 & [0.008, 0.017]\\
    Fitted lognormal   & 0.021 & 0.009 & 0.011 & [0.008, 0.015]\\
    \bottomrule
  \end{tabular}

  \vspace{0.6em}
  \raggedright
  The fitted service times give much smaller waits in this scenario.
  The intervals measure Monte Carlo error with the input models fixed.

  \vspace{0.4em}
  \textbf{Queue assumptions.} We count each entry to the automated phone
  menu as a request for an agent, even though some callers leave before
  reaching an agent. In the simulation, callers wait until served.
\end{frame}

\begin{frame}{Where would one extra agent make more difference?}
  The Week 3 reference demand is 39.85 calls per 15 minutes. Both fitted
  service models use a mean of 3.43 minutes.
  \[
    \frac{39.85}{15} \times 3.43
    \approx 9.11 \quad \text{agent minutes per minute}.
  \]

  \begin{itemize}
    \item With 16 agents, reference demand leaves about 6.9 agent minutes
      of spare capacity per minute.
    \item With 10 agents, it leaves about 0.9.
  \end{itemize}

  Use this average to choose a second experiment at \textbf{10 versus 11
  agents}. Simulated mornings can have quieter or busier periods than the
  reference demand.
\end{frame}

\begin{frame}{The extra agent helps more at lower staffing}
  \small
  For each staffing level, average the mean waiting time over the same
  200 simulated mornings.

  \vspace{0.6em}
  \centering
  \begin{tabular}{@{}lrr@{}}
    \toprule
    & \multicolumn{2}{c}{Mean waiting time (seconds)}\\
    \cmidrule(l){2-3}
    Agents & Fitted exponential & Fitted lognormal\\
    \midrule
    16 &   1.38 &   1.26\\
    17 &   0.62 &   0.57\\
    \addlinespace
    10 & 152.77 & 151.88\\
    11 &  64.45 &  66.08\\
    \bottomrule
  \end{tabular}

  \vspace{0.65em}
  \raggedright
  Going from 16 to 17 saves \textbf{less than one second} of mean wait.
  Going from 10 to 11 saves \textbf{about 86--88 seconds}.
\end{frame}

\begin{frame}{Lab: follow service times into the staffing decision}
  \begin{enumerate}
    \item Fit both service models and check their predictions on later dates.
    \item Measure how much work comes from long services.
    \item Calculate how much of the rise in long services the mix explains.
    \item Compare the gains from adding an agent at 10 and at 16 agents.
  \end{enumerate}

  \vspace{0.45em}
  Optional LLM audit. Predict what happens if you swap the group rates
  first and then the service mix. Check the prediction in your notebook.
\end{frame}

\begin{frame}{From individual inputs to the full bank case}
  \small
  \begin{tabular}{@{}lp{8.5cm}@{}}
    \textbf{Week 4} & Fit service durations and record where the candidates miss.\\[0.5em]
    \textbf{Week 5} & Estimate actual agent demand and caller patience. Add abandonment to the queue.\\[0.5em]
    \textbf{Week 6} & Check waiting and abandonment on later mornings. Compare one staffing change under the models that remain credible.\\[0.5em]
    \textbf{Week 7} & Present the complete workflow, evidence for the decision, and any remaining limits.\\
  \end{tabular}

  \vspace{0.65em}
  \href{https://github.com/NaitsabesMue/stochastic-modeling-simulation-ds}{Course repository and Python labs}
\end{frame}

\end{document}
